% ============================================================================
% 9 CONCLUSION AND OUTLOOK
% ============================================================================
\section{Conclusion and Outlook}
\label{sec:conclusion}

We built a reproducible route from raw text to evidence about an LLM system: a
configuration-driven evaluation framework, a document-scale SciREX pipeline, and
a human-in-the-loop annotation product. The central finding is not that one
prompting variant wins universally. It is that structural validity can coexist
with poor semantic accuracy, prompt definitions materially change behavior, and
quality must be measured on the target distribution.

The fresh SciREX result---$0.5738$ strict and $0.7030$ overlap $F_1$---supports
pre-annotation followed by review, not autonomous annotation. The 100-paper run
shows that sentence-local calls can reliably cover complete windows of more than
200 sentences with checkpointing, while making the scale visible: 5{,}868 calls and
2.91 million tokens for 100 papers. Reporting such operational evidence beside
accuracy is necessary for a defensible deployment decision.

\paragraph{Next steps.}
The immediate, feasible research steps are to instrument active annotation time
and editing actions, pilot a randomized manual-versus-pre-annotation study, and
use that pilot to determine the required sample size. The study should use at
least two domain experts, overlapping assignments, adjudication and unseen
target-domain tickets; it should measure final annotation quality and agreement
as well as time, rather than treating acceptance as correctness. In parallel,
the LLM, local nearest-neighbor and retrieval-augmented classifiers should be
evaluated on the same leakage-safe reference/development/test partition of
independently labeled internal tickets.

Next, improve Dataset/Metric definitions and test sentence groups where
cross-sentence context matters. The prepared 1{,}000-window SciREX study should
only be launched after actual provider pricing is configured; final train,
development and official-test results must be reported separately with
paper-clustered intervals.

Operational deployment is a larger engineering and governance phase, not a
small extension of the current Streamlit application. It requires organizational
authentication, project-scoped roles, concurrent assignment and version control,
a transactional database, background jobs, immutable audit logs, backup and
recovery, and tenant isolation. Before confidential tickets are processed, the
organization must also approve the complete data flow, model endpoint and
retention policy; enforce encryption, least privilege and managed secrets; and
verify provider data-use terms. These controls require the target organization's
infrastructure and security review and cannot be established by repository tests
alone. Fine-tuning is not a priority until a sufficiently large, stable and
consistently reviewed corpus exists. The Docker distribution and YAML use-case
layer make controlled research pilots reproducible without claiming readiness
for hundreds of concurrent users.
